% ---- score2 standard preamble (from: python3 score2_tools.py preamble) ----
\documentclass[10pt]{extarticle}
\usepackage[a4paper,margin=1.5cm]{geometry}
\usepackage[T1]{fontenc}
\usepackage{amsmath,amssymb}
\usepackage{graphicx}
\usepackage[export]{adjustbox}
\usepackage{multicol}
\usepackage{enumitem}
\usepackage{booktabs,array,tabularx}
\setlength{\columnsep}{0.7cm}
\setlength{\parindent}{0pt}
\setlength{\parskip}{3pt}
\pagestyle{empty}
\makeatletter
% figure registry (figdefs.tex is generated by the extract step)
\newcommand{\qfigdef}[4]{\@namedef{qf@file@#1}{#2}\@namedef{qf@w@#1}{#3}\@namedef{qf@wide@#1}{#4}}
\newcommand{\qf@check}[1]{\@ifundefined{qf@file@#1}{\PackageError{score2}{Unknown figure id #1}{}}{}}
\newcommand{\qf@img}[4][]{\includegraphics[width=#2\dimexpr\@nameuse{qf@w@#4}pt\relax,max width=#3,max height=0.45\textheight,#1]{\@nameuse{qf@file@#4}}}
% Full-width questions: finalize wraps every question that contains a wide figure in \qfullbegin...\qfullend,
% so the WHOLE question (stem, figures, options) is set across the page in reading order, then the two columns
% resume. The question's head - its number, stem and first figure (\qfullsplit marks the end of the head) -
% is kept on one page, so a number is never left at a page foot with its figure overleaf; the rest may break
% across pages normally, so no big gaps are left. Without a \qfullsplit, the whole question is kept together
% if it is under 45% of a page. Never write these yourself.
\newif\ifqf@full
\newif\ifqf@boxed
\newsavebox{\qf@box}
\newcommand{\qf@open}{\global\qf@fulltrue\par\global\qf@boxedtrue
  \setbox\qf@box\vbox\bgroup\hsize\textwidth\linewidth\textwidth}
\newcommand{\qf@emit}[1]{\par\egroup\global\qf@boxedfalse
  \ifdim\dimexpr\ht\qf@box+\dp\qf@box\relax<#1\textheight\noindent\usebox{\qf@box}\par\else\unvbox\qf@box\fi}
\newcommand{\qfullsplit}{\ifqf@boxed\qf@emit{0.85}\fi}
\newcommand{\qf@close}{\ifqf@boxed\qf@emit{0.45}\fi\global\qf@fullfalse}
\newcommand{\qfullbegin}{\end{multicols}\qf@open}
\newcommand{\qfullend}{\qf@close\begin{multicols}{2}}
\newcommand{\qfullstart}{\qf@open}
\newcommand{\qfulldone}{\qf@close}
% \qfig[scale]{id}: block figure, centred, directly where it is written (size set by the tool).
% Fallback if a wide figure is not inside a full-width question: break out of the columns in place - no floats.
\newcommand{\qfig}[2][1]{\qf@check{#2}%
  \ifqf@full\qf@blk{#1}{#2}\else
  \ifnum\@nameuse{qf@wide@#2}=1 \expandafter\@firstoftwo\else\expandafter\@secondoftwo\fi
  {\end{multicols}\begin{center}\qf@img{#1}{\textwidth}{#2}\end{center}\begin{multicols}{2}}%
  {\qf@blk{#1}{#2}}\fi}
\newcommand{\qf@blk}[2]{\par\smallskip{\centering\qf@img{#1}{\linewidth}{#2}\par}\smallskip}
% \qfigi[scale]{id}: inline figure (inside option lists / tables / minipages); never breaks columns.
\newcommand{\qfigi}[2][1]{\qf@check{#2}\qf@img[valign=c]{#1}{\linewidth}{#2}}
\makeatother
% structure macros - lint reads these to check completeness and figure placement
\newcommand{\papersec}[1]{\par\bigskip{\centering\Large\bfseries PAPER-#1\par}\medskip}
\newcommand{\partsec}[2]{\par\medskip\noindent{\large\bfseries PART-#1: #2}\par\smallskip}
\newcommand{\sect}[1]{\par\smallskip\noindent\textbf{#1}\par}
\newcommand{\qn}[1]{\par\smallskip\noindent\textbf{#1.}~}
\newcommand{\paraq}[2]{\par\smallskip\noindent\textbf{Paragraph for Questions #1 and #2}\par}
\newenvironment{qoptions}{\begin{itemize}[leftmargin=1.8em,itemsep=1pt,topsep=2pt,label={}]}{\end{itemize}}
\newcommand{\opt}[2]{\item[(#1)] #2}
% figure registry (generated by score2_tools.py finalize)
\qfigdef{p1_phy_s1_q01_1}{Score-II - Test 46 (26 Mar 2025)-img/p1_phy_s1_q01_1.png}{102.2}{0}
\qfigdef{p1_phy_s1_q02_1}{Score-II - Test 46 (26 Mar 2025)-img/p1_phy_s1_q02_1.png}{144.0}{0}
\qfigdef{p1_phy_s1_q02_2}{Score-II - Test 46 (26 Mar 2025)-img/p1_phy_s1_q02_2.png}{84.3}{0}
\qfigdef{p1_phy_s1_q02_3}{Score-II - Test 46 (26 Mar 2025)-img/p1_phy_s1_q02_3.png}{148.1}{0}
\qfigdef{p1_phy_s1_q02_4}{Score-II - Test 46 (26 Mar 2025)-img/p1_phy_s1_q02_4.png}{62.3}{0}
\qfigdef{p1_phy_s1_q03_1}{Score-II - Test 46 (26 Mar 2025)-img/p1_phy_s1_q03_1.png}{106.6}{0}
\qfigdef{p1_phy_s1_q04_1}{Score-II - Test 46 (26 Mar 2025)-img/p1_phy_s1_q04_1.png}{195.5}{0}
\qfigdef{p1_phy_s1_q05_1}{Score-II - Test 46 (26 Mar 2025)-img/p1_phy_s1_q05_1.png}{143.8}{0}
\qfigdef{p1_phy_s1_q06_1}{Score-II - Test 46 (26 Mar 2025)-img/p1_phy_s1_q06_1.png}{96.7}{0}
\qfigdef{p1_phy_s1_q07_1}{Score-II - Test 46 (26 Mar 2025)-img/p1_phy_s1_q07_1.png}{120.4}{0}
\qfigdef{p1_phy_s1_q09_1}{Score-II - Test 46 (26 Mar 2025)-img/p1_phy_s1_q09_1.png}{124.0}{0}
\qfigdef{p1_phy_s1_q10_1}{Score-II - Test 46 (26 Mar 2025)-img/p1_phy_s1_q10_1.png}{109.6}{0}
\qfigdef{p1_phy_s2_q01_1}{Score-II - Test 46 (26 Mar 2025)-img/p1_phy_s2_q01_1.png}{95.1}{0}
\qfigdef{p1_phy_s2_q02_1}{Score-II - Test 46 (26 Mar 2025)-img/p1_phy_s2_q02_1.png}{148.4}{0}
\qfigdef{p1_phy_s2_q03_1}{Score-II - Test 46 (26 Mar 2025)-img/p1_phy_s2_q03_1.png}{109.9}{0}
\qfigdef{p1_phy_s2_q06_1}{Score-II - Test 46 (26 Mar 2025)-img/p1_phy_s2_q06_1.png}{91.6}{0}
\qfigdef{p1_phy_s2_q07_1}{Score-II - Test 46 (26 Mar 2025)-img/p1_phy_s2_q07_1.png}{91.8}{0}
\qfigdef{p1_chem_s1_q01_A}{Score-II - Test 46 (26 Mar 2025)-img/p1_chem_s1_q01_A.png}{276.3}{1}
\qfigdef{p1_chem_s1_q01_C}{Score-II - Test 46 (26 Mar 2025)-img/p1_chem_s1_q01_C.png}{227.3}{0}
\qfigdef{p1_chem_s1_q01_D}{Score-II - Test 46 (26 Mar 2025)-img/p1_chem_s1_q01_D.png}{75.4}{0}
\qfigdef{p1_chem_s1_q03_1}{Score-II - Test 46 (26 Mar 2025)-img/p1_chem_s1_q03_1.png}{123.7}{0}
\qfigdef{p1_chem_s1_q04_1}{Score-II - Test 46 (26 Mar 2025)-img/p1_chem_s1_q04_1.png}{162.0}{0}
\qfigdef{p1_chem_s1_q07_1}{Score-II - Test 46 (26 Mar 2025)-img/p1_chem_s1_q07_1.png}{165.4}{0}
\qfigdef{p1_chem_s1_q07_2}{Score-II - Test 46 (26 Mar 2025)-img/p1_chem_s1_q07_2.png}{258.2}{1}
\qfigdef{p1_chem_s1_q09_1}{Score-II - Test 46 (26 Mar 2025)-img/p1_chem_s1_q09_1.png}{267.2}{1}
\qfigdef{p1_chem_s1_q09_2}{Score-II - Test 46 (26 Mar 2025)-img/p1_chem_s1_q09_2.png}{129.9}{0}
\qfigdef{p1_chem_s1_q09_3}{Score-II - Test 46 (26 Mar 2025)-img/p1_chem_s1_q09_3.png}{119.9}{0}
\qfigdef{p1_chem_s1_q09_4}{Score-II - Test 46 (26 Mar 2025)-img/p1_chem_s1_q09_4.png}{129.9}{0}
\qfigdef{p1_chem_s1_q09_5}{Score-II - Test 46 (26 Mar 2025)-img/p1_chem_s1_q09_5.png}{124.0}{0}
\qfigdef{p1_chem_s2_q03_1}{Score-II - Test 46 (26 Mar 2025)-img/p1_chem_s2_q03_1.png}{322.1}{1}
\begin{document}
\begin{center}{\Large\bfseries JEE (Main+Advanced) Score-II --- Test 46 (26 Mar 2025)}\end{center}
\begin{multicols}{2}

\partsec{1}{PHYSICS}
\sect{Section-I}

\qn{1} A geometric progression of very large number of identical springs of spring constant k is pulled with a force F. Assume bars \& springs are massless. Last row is connected to rigid support. In equilibrium choose the correct option(s)
\qfig{p1_phy_s1_q01_1}
\begin{qoptions}
\opt{A}{Free end is displaced by amount $\dfrac{2F}{k}$.}
\opt{B}{Free end is displaced by amount $\dfrac{F}{k}$.}
\opt{C}{Displacement of bar PQ is $\dfrac{F}{k}$.}
\opt{D}{Displacement of bar PQ is $\dfrac{F}{3k}$.}
\end{qoptions}

\qn{2} The figure below shows a practical situation of fluid flow where Bernoulli's principle holds good.
\qfig{p1_phy_s1_q02_1}
If $\dfrac{p}{\rho g}$ is pressure head, $\dfrac{v^2}{2g}$ is kinetic head, h is potential head, then variation of energy heads with respect to 'x' (coordinate) are as shown
\qfig{p1_phy_s1_q02_2}
\qfig{p1_phy_s1_q02_3}
\qfig{p1_phy_s1_q02_4}
Choose correct option(s)
\begin{qoptions}
\opt{A}{Kinetic head v/s x will be best represented by curve (4)}
\opt{B}{Potential head v/s x will be best represented by curve (2)}
\opt{C}{Pressure head v/s x will be best represented by curve (3)}
\opt{D}{$\left(\dfrac{p}{\rho g}\right)_P - \left(\dfrac{p}{\rho g}\right)_Q = 70$ m}
\end{qoptions}

\qn{3} A bridge for measuring the inductance $L_1$ \& resistance $R_1$ of a coil is shown in figure. $R_2$ \& $R_3$ are variable resistors \& $C_3$ \& $C_4$ are fixed capacitor. For bridge to be balanced, choose the correct option(s)
\qfig{p1_phy_s1_q03_1}
\begin{qoptions}
\opt{A}{$R_1 = \dfrac{R_2 C_4}{C_3}$}
\opt{B}{$L_1 = C_4 R_3 R_2$}
\opt{C}{Conditions for the bridge to be balanced are independent of the frequency of the alternating voltage source V.}
\opt{D}{Conditions for the bridge to be balanced are dependent of the frequency of the alternating voltage source V.}
\end{qoptions}

\qn{4} Following graphs represent the frequency dependence of the amplitude of displacements and frequency dependence of the velocity amplitude in forced oscillation on the frequency of the driving force. Choose the correct option(s)
\qfig{p1_phy_s1_q04_1}
\begin{qoptions}
\opt{A}{Figure (a) correspond to frequency dependence of the displacement amplitudes}
\opt{B}{Figure (b) correspond to frequency dependence of the displacement amplitudes}
\opt{C}{Figure (b) correspond to the frequency dependence of the velocity amplitudes.}
\opt{D}{Figure (a) correspond to the frequency dependence of the velocity amplitudes.}
\end{qoptions}

\qn{5} Each wire has Resistance = R. Choose the correct option(s)
\qfig{p1_phy_s1_q05_1}
\begin{qoptions}
\opt{A}{Equivalent resistance between P \& R is $\dfrac{8R}{15}$}
\opt{B}{Equivalent resistance between Q \& T is $\dfrac{3R}{5}$}
\opt{C}{Equivalent resistance between P \& R is $\dfrac{15R}{8}$}
\opt{D}{Equivalent resistance between Q \& T is $\dfrac{4R}{5}$}
\end{qoptions}

\qn{6} A small body is placed on the top of smooth hemisphere of radius R. The hemisphere is imparted a constant acceleration $a_0=g$ in horizontal direction as shown.
\qfig{p1_phy_s1_q06_1}
\begin{qoptions}
\opt{A}{The velocity of the body relative to sphere at the moment of break off is $V_0=\sqrt{\dfrac{4gR}{3}}$}
\opt{B}{The velocity of the body relative to sphere at the moment of break off is $V_0=\sqrt{\dfrac{2gR}{3}}$}
\opt{C}{The angle $\theta_0$ between the vertical \& the radius vector drawn from the centre of sphere to break off point is given by $\cos\theta_0 = \dfrac{2+\sqrt{14}}{6}$.}
\opt{D}{The angle $\theta_0$ between the vertical \& the radius vector drawn from centre of the sphere to the break off point is given by $\cos\theta_0=\dfrac{2}{3}$}
\end{qoptions}

\qn{7}
\qfig{p1_phy_s1_q07_1}
In a photoelectric effect setup, we have 4 metals which can be used for emitter or collector. (hc = 1240 eVnm)
\begin{itemize}[leftmargin=1.8em,itemsep=0pt,topsep=2pt]
\item[(1)] Aluminium $\phi=3.7$ eV
\item[(2)] Platinum $\phi=5.3$ eV
\item[(3)] Silver $\phi=4.3$ eV
\item[(4)] Cesium $\phi=1.9$ eV
\end{itemize}
We use monochromatic light $\lambda=310$ nm. In following options, choose the cases in which no photocurrent is detected. V is the voltage of the battery used with polarity as shown.
\begin{qoptions}
\opt{A}{Emitter : Al, collector Cs, V = 1V}
\opt{B}{Emitter : Cs, collector Pt, V = 2V}
\opt{C}{Emitter : Ag, collector Cs, V = 1V}
\opt{D}{Emitter : Pt, collector Cs, V = $-2$V}
\end{qoptions}

\qn{8} The following equations represent transverse wave
\begin{center}
\resizebox{\linewidth}{!}{$Z_1 = a\cos(kx-\omega t)\qquad Z_2 = a\cos(kx+\omega t)\qquad Z_3 = a\cos(ky-\omega t)$}
\end{center}
Choose correct option(s)
\begin{qoptions}
\opt{A}{For the combination of $Z_1$ \& $Z_3$, the wavefronts are represented by the planes $x+y=$ constant}
\opt{B}{$Z_1+Z_3$ represents a wave going along the direction making an angle $45^\circ$ with the positive x \& positive y-direction}
\opt{C}{$Z_1+Z_2$ represents a standing wave.}
\opt{D}{For the combination of $Z_2$ \& $Z_3$, the wavefronts are represented by the planes $y-x=$ constant}
\end{qoptions}

\qn{9} A nonconducting solid sphere A and a nonconducting hollow sphere B are placed on rough insulated surface at a very large distance as shown. Spheres start moving under the mutual force of attraction and collide elastically. (Assume no slipping any where)
\qfig{p1_phy_s1_q09_1}
\begin{qoptions}
\opt{A}{Velocity of A just before collision is $\sqrt{\dfrac{126}{335}\dfrac{KQ^2}{mR}}$}
\opt{B}{Velocity of A just before collision is $\dfrac{25}{42}\sqrt{\dfrac{126}{335}\dfrac{KQ^2}{mR}}$}
\opt{C}{Velocity of B just before collision is $\sqrt{\dfrac{126}{335}\dfrac{KQ^2}{mR}}$}
\opt{D}{Velocity of B just before collision is $\dfrac{25}{42}\sqrt{\dfrac{126}{335}\dfrac{KQ^2}{mR}}$}
\end{qoptions}

\qn{10} A non conducting cylinder is divided in two equal parts by a non conducting movable piston. One part contains monoatomic gas and other part contains diatomic gas, each of 1 mole. Now heat is supplied to the left part such that pressure of the right part becomes $\dfrac{243P_0}{32}$. Then :-
\qfig{p1_phy_s1_q10_1}
\begin{qoptions}
\opt{A}{Final temperature of gas in right chamber will be $\dfrac{27}{8}T_0$}
\opt{B}{Final temperature of gas in left chamber will be $\dfrac{207}{16}T_0$}
\opt{C}{Work done by gas in the right chamber will be $\dfrac{-15}{8}P_0V_0$}
\opt{D}{Heat supplied by heater will be $\dfrac{1015}{32}P_0V_0$}
\end{qoptions}

\sect{Section-II}

\qn{1} In the potentiometer circuit shown, the battery has emf V = 10 V and internal resistance of 5 $\Omega$. R is a resistance box whose resistance can be any integral value from 1 $\Omega$ to 100 $\Omega$. Resistance of potentiometer wire is 50 $\Omega$ and length is 5 m. The emf E of unknown battery can be anywhere between 5 V to 7 V. For most accurate measurement of emf E, what should be value of R (in ohm)?
\qfig{p1_phy_s2_q01_1}

\qn{2} A rectangular tank of height h filled with water is placed near the bottom of a plane mirror inclined at an angle $45^\circ$ with the horizontal. At height 10 m from bottom a hole is made such that stream coming out from hole strikes the plane mirror normally. Then at the time of hitting, speed of image of water droplet hitting the plane mirror, formed in plane mirror (in m/s) is (Take $g=10$ m/s$^2$) :-
\qfig{p1_phy_s2_q02_1}

\qn{3} Consider a block suspended with the light rod being horizontal in equilibrium as shown. All springs are light. Neglect friction everywhere. The time period for small oscillation of the block is $2\pi\sqrt{\dfrac{\alpha M}{5K}}$. Find $\alpha$.
\qfig{p1_phy_s2_q03_1}

\qn{4} Atoms of type A fly along the axis of a cylindrical pipe of radius R and collide with practically motionless atoms of type B. The kinetic energy of atoms A is equal to the threshold energy, so that a central impact forms a molecule AB, which then moves with velocity v. In the case of an off-centre impact, the reaction does not occur, i.e. the atoms collide elastically. In what minimum time (in sec) after the collision will atoms of type B be able to get from the axis of the cylinder to the wall of the cylinder? (R = 10 cm and v = 40 cm/s)

\qn{5} The density of sphere is directly proportional to depth below the surface of the sphere. If the maximum value of acceleration due to gravity is N times of the surface value, then the value of N is

\qn{6} A rod of resistance R, length $\ell$ and mass m is placed on two smooth parallel rails which are resistanceless. The rails are connected by capacitor of capacitance C. A uniform magnetic field B is present in vertical direction. The rod is pulled at it's center by a constant force F. If $m=B^2\ell^2 C$, RC = 2 sec, F = 7.29 N, B = 0.1 T and $\ell=2$ m, what is current (in A) in rod at t = 2 sec. (e = 2.7)
\qfig{p1_phy_s2_q06_1}

\qn{7} A small block of mass m \& charge q is released on a long rough inclined plane as shown. In space, there is a uniform magnetic field of 0.5 T away from the reader. Find the maximum velocity (in m/s) of the mass during the motion. (B = 0.5 T, q = 4$\mu$C, m = 8 mg, g = 10 m/s$^2$)
\qfig{p1_phy_s2_q07_1}

\qn{8} A well lagged wire of length L = 2 m of cross sectional area A = 4 m$^2$ has it's ends maintained at temperatures $T_1=40^\circ$C \& $T_2=10^\circ$C. The thermal conductivity of the wire is given by $K=(3+6T)\,\dfrac{\mathrm{W}}{\mathrm{m}^\circ\mathrm{C}}$, where T is the temperature (in $^\circ$C). The rate of flow of heat (in J/s) along the wire is:

\qfullbegin
\partsec{2}{CHEMISTRY}
\sect{Section-I}

\qn{1} Which of the following statement(s) is/are \textbf{incorrect} ?
\begin{qoptions}
\opt{A}{\qfigi{p1_chem_s1_q01_A}}
\opt{B}{1,3,5-trihydroxybenzene reacts with NH$_2$OH but 1,2,4-trihydroxybenzene and 1,2,3-trihydroxybenzene do not react with NH$_2$OH.}
\opt{C}{\qfigi{p1_chem_s1_q01_C}}
\opt{D}{\qfigi{p1_chem_s1_q01_D} number of Iodoform molecules produced per molecule of A is 3.}
\end{qoptions}
\qfullend

\qn{2} Which epimer or epimers of D-glucose gives meso aldaric acid as product with conc. HNO$_3$.
\begin{qoptions}
\opt{A}{C-2 Epimer}
\opt{B}{C-3 Epimer}
\opt{C}{C-4 Epimer}
\opt{D}{C-5 Epimer}
\end{qoptions}

\qn{3} Identify the correct statement(s) from the following, at constant temperature :
\qfig{p1_chem_s1_q03_1}
A = Standard free energy of 1 mole of N$_2$O$_4$; C = standard free energy of 2 mole of NO$_2$; B = Equilibrium point
\begin{qoptions}
\opt{A}{100\% conversion of dimer to monomer is spontaneous.}
\opt{B}{At point B, free energy of system becomes constant with time.}
\opt{C}{In the region between A and B, reaction has tendency to move in forward direction.}
\opt{D}{In the region between B and C, $\Delta G$ for the reaction $\mathrm{N_2O_4(g)\rightleftharpoons 2NO_2(g)}$ is positive.}
\end{qoptions}

\qn{4} Which of the following statements regarding A is/are incorrect ?
\qfig{p1_chem_s1_q04_1}
\begin{qoptions}
\opt{A}{Benzenoid form of anion of A is pinkish red in colour.}
\opt{B}{Quinonoid form of anion of A is yellow in colour.}
\opt{C}{Benzenoid form of anion accepts a proton and gets converted into a more deep coloured quinonoid form.}
\opt{D}{A is used as indicator in weak base v/s strong acid titration.}
\end{qoptions}

\qn{5} In CaF$_2$ type crystal, which of following statement(s) is/are true?
\begin{qoptions}
\opt{A}{Anion is inside a cubical void made by cation.}
\opt{B}{For each cation, there will be 12 cations which are its second nearest neighbours.}
\opt{C}{For each cation, there will be 16 third nearest as anionic neighbours.}
\opt{D}{The third nearest neighbour of cation is at a distance of $\dfrac{\sqrt{15}}{4}a$.}
\end{qoptions}

\qn{6} $\left[\mathrm{Co(H_2O)_6}\right]^{2+} + \mathrm{Cl}^- \rightleftharpoons \text{"X"} + 6\mathrm{H_2O}$

Correct statement(s) about the reaction is/are
\begin{qoptions}
\opt{A}{X is blue coloured compound with octahedral structure.}
\opt{B}{Increasing the temperature favours forward reaction.}
\opt{C}{Addition of Ca$^{2+}$ to the equilibrium mixture favours forward reaction.}
\opt{D}{Addition of Zn$^{2+}$ ions to the equilibrium mixture favours forward reaction.}
\end{qoptions}

\qfullbegin
\qn{7}
\qfig{p1_chem_s1_q07_1}
\qfullsplit
Note: C and D are oxide of same element.
\qfig{p1_chem_s1_q07_2}
Which of the following option(s) is/are correct ?
\begin{qoptions}
\opt{A}{Hybridization of metal in complex \textcircled{M} is dsp$^2$.}
\opt{B}{Basicity of compound \textcircled{K} is two.}
\opt{C}{Compound \textcircled{C} is H$_2$SO$_4$.}
\opt{D}{Oxidation state of central metal ion in compound \textcircled{F} is +3.}
\end{qoptions}
\qfullend

\qn{8} $\alpha$ - maltose (C$_{12}$H$_{22}$O$_{11}$) can be hydrolysed to glucose (C$_6$H$_{12}$O$_6$) according to the following reaction:
\[\mathrm{C_{12}H_{22}O_{11}(aq) + H_2O(\ell) \rightarrow 2\,C_6H_{12}O_6(aq)}\]
Given: Standard enthalpy of formation of C$_{12}$H$_{22}$O$_{11}$(aq) = $-2238$ kJ/mol; Standard enthalpy of formation of H$_2$O($\ell$) = $-285$ kJ/mol; Standard enthalpy of formation of C$_6$H$_{12}$O$_6$(aq) = $-1263$ kJ/mol
\begin{center}
\begin{tabular}{|l|c|c|c|}
\hline
Time (min) & 0 & 50 & 100 \\
\hline
Conc.\ of $\alpha$-maltose (M) & 4.0 & 1.0 & 0.25 \\
\hline
\end{tabular}
\end{center}
which of the following statement(s) is/are true?
\begin{qoptions}
\opt{A}{The hydrolysis of $\alpha$-maltose is exothermic.}
\opt{B}{Heat liberated in combustion of 1 mole of $\alpha$-maltose(s) must be greater than the heat liberated in combustion of 2 moles of glucose(s).}
\opt{C}{Increasing temperature will increase the degree of hydrolysis of $\alpha$-maltose.}
\opt{D}{The hydrolysis of $\alpha$-maltose follows $1^{\mathrm{st}}$ order kinetics.}
\end{qoptions}

\qfullbegin
\qn{9} Vanillin (L), C$_8$H$_8$O$_3$, is isolated from vanilla beans. It give an intense colour with FeCl$_3$ and positive Tollen's test. It is not steam-distilled and does not react with HCl. It goes through following steps :
\qfig{p1_chem_s1_q09_1}
\qfullsplit
Which of the following is correct statement about various products?
\begin{qoptions}
\opt{A}{L is \qfigi{p1_chem_s1_q09_2}}
\opt{B}{L is \qfigi{p1_chem_s1_q09_3}}
\opt{C}{L is \qfigi{p1_chem_s1_q09_4}}
\opt{D}{L is \qfigi{p1_chem_s1_q09_5}}
\end{qoptions}
\qfullend

\qn{10} White phosphorous reacts with NaOH to produce gas (X) and the compound (Y).

'X' + HCl $\longrightarrow$ 'Z'

What are the changes take place when 'X' is converted to Z.
\begin{qoptions}
\opt{A}{State of hybridisation of central atom.}
\opt{B}{\% s-character of bonds associated to central atom.}
\opt{C}{Oxidation state of central atom}
\opt{D}{Number of lone pair(s) on central atom.}
\end{qoptions}

\sect{Section-II}

\qn{1} If number of molecules per unit volume reduces by 20\% and radius of molecule increased by 20\%, the ratio of new mean free path to the old mean free path is x : 1. The value of 'x' is.

\qn{2} How many of the following statement(s) is/are true with respect to preparation of sodium hydroxide by electrolysis of brine solution with graphite anode ?
\begin{enumerate}[label=(\roman*),leftmargin=2.2em,itemsep=1pt,topsep=2pt]
\item H$_2$ is released at cathode, if we use Fe at cathode.
\item With Hg cathode, Na$^+$ is deposited instead of H$^+$ ion.
\item Standard reduction potential data suggests that it is easier to deposit H$^+$ than Na$^+$ ion.
\item Hg has a low hydrogen overvoltage while sodium has a very high overvoltage on Hg.
\item Cl$_2$ is released at anode due to high overvoltage of O$_2$.
\item It is advised not to use an open cell for the above mentioned electrolysis.
\end{enumerate}

\qfullbegin
\qn{3}
\qfig{p1_chem_s2_q03_1}
\qfullsplit
Calculate the gm of aromatic product (E), if 1865 gm of reactant (X) is used for above reaction sequence.
[\textbf{Given, atomic weight :} H = 1, C = 12, O = 16, F = 19, Cl = 35.5, Br = 80]
\qfullend

\qn{4} Number of moles of periodic acid required to oxidise one mole of glucose is X and number of moles of formic acid formed when one mole of fructose get oxidised by enough periodic acid is Y. Find the value of (X $-$ Y).

\qn{5} How many of the following ions will form white precipitate with K$_4$[Fe(CN)$_6$] (before aerial oxidation)?
\[\mathrm{Co^{2+},\ Ca^{2+},\ Fe^{2+},\ Fe^{3+},\ Mn^{2+},\ Ag^+,\ Zn^{2+}}\]

\qn{6} How many of the following statements are correct?
\begin{enumerate}[label=(\arabic*),leftmargin=2.2em,itemsep=1pt,topsep=2pt]
\item Micellization is entropy driven process.
\item A catalyst has no effect on $\Delta G$, $K_{eq}$ and $\Delta H$ of a reaction.
\item On increasing non-polar chain length of R$-$COO$^-$Na$^+$, CMC increases.
\item Heamoglobin is a positive sol.
\item Charcoal is negative sol.
\end{enumerate}

\qn{7} For the reaction : $4A(g) + 2B(g) \rightarrow 3C(g)$; $\Delta_r S^\circ = -40$ cal/K-mol at 300 K. Calculate $\Delta_r S^\circ$ for the reaction at 400 K, in cal K$^{-1}$ mol$^{-1}$.

(Given : $C_{P,m}$ [/Cal K$^{-1}$ mol$^{-1}$] : A(g) = 3.0, B(g) = 5.0; C(g) = 7.0; $\ln2=0.7$; $\ln3=1.1$)

\qn{8} NH$_3$ is not obtained as product from how many of the following reactions.
\begin{center}
\begin{tabular}{@{}ll@{}}
(a) $(\mathrm{NH_4})_2\mathrm{HPO_4} \xrightarrow{\Delta}$ & (d) $\mathrm{B_3N_3H_6} + \mathrm{H_2O} \longrightarrow$ \\[4pt]
(b) $\mathrm{CaCN_2} + \mathrm{H_2O} \longrightarrow$ & (e) $\mathrm{Ba(N_3)_2} \xrightarrow{\Delta}$ \\[4pt]
(c) $\mathrm{NH_4ClO_4} \xrightarrow{\Delta}$ & (f) $\mathrm{Ba_3N_2} + \mathrm{H_2O} \longrightarrow$
\end{tabular}
\end{center}

\partsec{3}{MATHEMATICS}
\sect{Section-I}

\qn{1} Which of the following is/are true
\begin{qoptions}
\opt{A}{$\displaystyle\int_0^{\pi/2}\dfrac{\sin(7x)}{\sin x}\,dx = \dfrac{\pi}{2}$}
\opt{B}{$\displaystyle\int_0^{\pi/2}\dfrac{\sin(10x)}{\sin x}\,dx = \dfrac{526}{315}$}
\opt{C}{$\displaystyle\int_0^{\pi}\dfrac{\sin(9x)}{\sin x}\,dx = \pi$}
\opt{D}{$\displaystyle\int_0^{1} x^4(1-x)^4\,dx = \dfrac{1}{630}$}
\end{qoptions}

\qn{2} Let A and B are two non-zero $n\times n$ matrices then the value of $\dfrac{\mathrm{Tr}\left[(AB-BA)^n\right]}{\det(AB-BA)}$ equals (where Tr(x) and det(x) represents Trace and determinant value of matrix-x) (where $\det(AB-BA)\neq0$)
\begin{qoptions}
\opt{A}{$-2$ If n = 2}
\opt{B}{3 If n = 3}
\opt{C}{$\dfrac{1}{3}$ If n = 3}
\opt{D}{$\dfrac{1}{2}$ If n = 2}
\end{qoptions}

\qn{3} Consider a tetrahedron is formed by planes $y+z=0$, $z+x=0$, $x+y=0$ and $x+y+z=1$ which of the following is/are true
\begin{qoptions}
\opt{A}{volume of tetrahedron is $\dfrac{4}{3}$}
\opt{B}{centroid of tetrahedron is $\left(\dfrac{1}{4},\dfrac{1}{4},\dfrac{1}{4}\right)$}
\opt{C}{Shortest distance between two of non intersecting edges of tetrahedron is $\sqrt{\dfrac{2}{3}}$}
\opt{D}{One of the edge of tetrahedron is $\dfrac{x-1}{0}=\dfrac{y+1}{2}=\dfrac{z-1}{-2}$}
\end{qoptions}

\qn{4} A curve $y=f(x)$ passes through point P(1,1), the normal to the curve at P is $a(y-1)+(x-1)=0$. If slope of tangent at any point on the curve is proportional to ordinate of that point then correct option(s) is/are -
\begin{qoptions}
\opt{A}{Equation of the curve is $y=e^{a(x-1)}$ (whre $a\in\mathbb{R}$)}
\opt{B}{Area bounded by y-axis, curve and normal to curve at P is $\left|1+\dfrac{e^{-a}}{a}-\dfrac{1}{2a}\right|$}
\opt{C}{Area bounded by f(x), x-axis in $x\in(0,1)$ for some value of 'a' is $\left(1-\dfrac{1}{e}\right)$}
\opt{D}{Area bounded by y-axis, curve and normal to curve at P is $\left|\displaystyle\int_0^1\left(1+\dfrac{1}{a}-\dfrac{x}{a}-e^{a(x-1)}\right)dx\right|$}
\end{qoptions}

\qn{5} Let the mirror image of the curve $\dfrac{(x-1)^2}{4}+\dfrac{(y+2)^2}{9}=1$, with respect to the line $x+y-4=0$ is shifted 1 unit right and 2 unit up and the resulting curve is C which of the following regarding the curve C is/are correct?
\begin{qoptions}
\opt{A}{Centre of C is (7,5)}
\opt{B}{Latus rectum of C is 9/2}
\opt{C}{The curve C also passes through $\left(7+\dfrac{3}{\sqrt2},\,5+\sqrt2\right)$}
\opt{D}{The point (4, 5) lies inside C.}
\end{qoptions}

\qn{6} Consider a curve mirror $y=f(x)$ passing through (8, 6) having property that all rays emanating from (0, 0) after getting reflected from mirror becomes parallel to x-axis then correct option is/are
\begin{qoptions}
\opt{A}{Length of latus rectum of f(x) can be 4}
\opt{B}{Length of latus rectum of f(x) can be 36}
\opt{C}{Directrix of f(x) can be $x=-2$}
\opt{D}{Directrix of f(x) can be $x=16$}
\end{qoptions}

\qn{7} Consider the set A = $\{1,2,3,4,5\}$.
Let the number of functions defined from A to A for which $f(f(i))=i$ for at least one $i\in A$ is N then
\begin{qoptions}
\opt{A}{N has four divisors}
\opt{B}{N is odd}
\opt{C}{N is even}
\opt{D}{N is of the form $4\lambda+1$ (where $\lambda\in\mathbb{N}$)}
\end{qoptions}

\qn{8} Which of the following statement is/are true?
\begin{qoptions}
\opt{A}{If $f(g(x))$ is injective then $g(x)$ must be injective}
\opt{B}{If $f(g(x))$ is surjective then $f(x)$ must be surjective}
\opt{C}{If $f(g(x))$ is injective $g(x)$ may not be injective}
\opt{D}{If $f(g(x))$ is surjective $f(x)$ may not be surjective}
\end{qoptions}

\qn{9} Let x, y and z are three integers satisfying the system of simultaneous equations
$xy+z=100$, $yz+x=87$ and $zx+y=60$.
Which of the following option(s) is/are correct?
\begin{qoptions}
\opt{A}{$x+y+z=-29$}
\opt{B}{$x+y+z=-28$}
\opt{C}{$xy+yz+zx=276$}
\opt{D}{$xy+yz+zx=275$}
\end{qoptions}

\qn{10} Let 'f' be a polynomial function such that $f(f(x))-8x^m=\dfrac{80}{f''(x)}-64x^2+100\ \forall x\in\mathbb{R}$ and $m\in\mathbb{N}$.

If $f(|x|)$ is differentiable for all $x\in\mathbb{R}$ then choose the correct option(s).
\begin{qoptions}
\opt{A}{m = 4}
\opt{B}{m = 2}
\opt{C}{$\left[\dfrac{f(3)}{m}\right]=2$, (where [.] denotes greatest integer function)}
\opt{D}{f(x) is a polynomial of odd degree.}
\end{qoptions}

\sect{Section-II}

\qn{1} Number of permutations of 1, 2, 3, 4, 5, 6 such that for each k with $1\le k\le5$, at least one of the first k terms of the permutation is greater than k.

\qn{2} Let f(x,y,z) be a Homogenous polynomial of degree-4 in real variables x, y, z satisfying the following conditions
(a) $f(1,2,3)=1/10$
(b) $f(x,y,z)=f(z,x,y)=f(y,z,x)$ for all $x,y,z\in\mathbb{R}$
(c) $f(x,x,z)=0$ for all $x,z\in\mathbb{R}$

Value of $f(1,4,9)$ equals

\qn{3} Let $T_n = 1+2+3+\cdots+n$ and $M_n = \dfrac{T_2}{T_2-1}\times\dfrac{T_3}{T_3-1}\times\dfrac{T_4}{T_4-1}\times\cdots\times\dfrac{T_n}{T_n-1}$ for $n=2,3,4,\ldots$

Value of $\dfrac{13}{25}M_{50}$ equals

\qn{4} $\dfrac{\tan^6\left(\dfrac{\pi}{9}\right)+27\tan^2\left(\dfrac{\pi}{9}\right)}{1+11\tan^4\left(\dfrac{\pi}{9}\right)} = \sqrt{P}-1$ (where $P\in\mathbb{N}$)

then the remainder when P is divided by 7 is\underline{\hspace{1.5cm}}

\qn{5} Consider the function
\[f(x)=\max\left\{x^3-\dfrac12,\,-\sqrt{16-x^2},\,\dfrac{x}{|x|}\right\},\ x\in[-4,4]\]
m = number of points where the function is discontinuous
n = number of points where the function is continuous but not differentiable
find (m + n)

\qn{6} In a circle 15 persons are seated and it is known that each person knows the two immediate neighbor (one on the left and one on the right) only. If 5 persons are selected at random, then the probability that the selected persons are unknown to each other is $\dfrac{p}{q}$, where p, q are coprime then sum of digits of p is\underline{\hspace{1.5cm}}

\qn{7} If $(1+x)^n = P_0+P_1x+P_2x^2+P_3x^3+\cdots,\ (n\in\mathbb{N})$

$P_0-P_2+P_4-\cdots=\alpha$, $P_1-P_3+P_5-\cdots=\beta$ then $\displaystyle\sum_{n=1}^{\infty}\dfrac{n}{\alpha^2+\beta^2}$ is

\qn{8} If $A(\vec a), B(\vec b), C(\vec c), D(\vec d)$ are position vector of a cyclic quadrilateral ABCD,
\[\alpha=\dfrac{\left|(\vec a\times\vec b)+(\vec b\times\vec d)+(\vec d\times\vec a)\right|}{\left|(\vec b-\vec a)\cdot(\vec d-\vec a)\right|}\]
\[\beta=\dfrac{\left|(\vec b\times\vec c)+(\vec c\times\vec d)+(\vec d\times\vec b)\right|}{\left|(\vec b-\vec c)\cdot(\vec d-\vec c)\right|}\]
then $\left|\dfrac{3\alpha}{2\beta}\right|$ is

\end{multicols}
\end{document}
